\centering
\begin{tabular}{m{1cm}  m{1.5cm} m{2.2cm}  m{8cm} }
\hline
 \textbf{Date}           & \textbf{State} &   \textbf{Affected group} & \textbf{Description}   \\ \hline
 April 2016 & Lower-Austria & Subsidiary-protected refugees & Eligibility for minimum income support was revoked. A single
adult who received 838 euros, while on minimum income support, now only gets 365 euros if living on their own, or has the option to stay in a refugee camp. \\\hline
July 2016 & Upper Austria & Both groups & A new minimum income standard was introduced for subsidiary-protected and convention refugees, who came to Austria after November 15th, 2015. The new minimum income standard was 522 euros. In November 2018, the reform was revoked by the European Court of Justice. \\\hline
January 2017 & Lower Austria  & Convention refugees & A new minimum income standard for people who have spent less than 5 out of the last 6 years in Austria was introduced. For a single adult, the new minimum standard was 572,50 euros. For families and people living in a shared flat, the state introduced an additional cap of 1,500 euros for the total a household could claim, irrespective of the number of people living there. In March 2018, the reform was revoked by the Constitutional Court. \\\hline
April 2017  & Burgenland & Convention refugees & A new minimum income standard for people who have spent less than 5 out of the last 6 years in Austria was introduced. For a single adult, the new minimum standard was 585 euros. For families and people living in a shared flat, the state introduced an additional cap of 1,500 euros for the total a household could claim, irrespective of the number of people living there. In March 2018, the reform was revoked by the Constitutional Court. \\
 
\hline


\end{tabular}